\begin{abstract}
High-resolution national hydrography and environmental datasets make detailed SWAT\textsuperscript{+} model construction possible across broad spatial extents, but reproducible model assembly remains difficult because routing topology, attributes, and export structure can vary across analysts and study sites.
This paper evaluates SWATGenX as a reference implementation of an automated NHDPlus High Resolution (NHDPlus HR) workflow that converts U.S.\ national vectors and ancillary layers into SWAT\textsuperscript{+}-compatible project structures with explicit preprocessing rules and quality-assurance checks.
\textbf{Primary evidence} documents implementation auditability and quantified hydrography preprocessing trade-offs, including a drainage-area fidelity audit that compares executable SWAT\textsuperscript{+} channel areas with original NHDPlus HR cumulative drainage at gaged channels (Objectives~1--2).
\textbf{Secondary evidence} reports structural scaling, model-assembly cost, and downstream workflow readiness---initialization, split-sample calibration and verification, and Morris sensitivity on two controlled gage basins (Objectives~3--4).
\textbf{Supporting evidence} reports measured SWAT\textsuperscript{+} simulation cost on one documented host/build configuration (Objective~5).
Hydrologic skill metrics are workflow diagnostics on selected basins, not a national validation study.
SWATGenX demonstrates that high-resolution SWAT\textsuperscript{+} project generation from NHDPlus HR and national datasets can be automated into executable, auditable model packages with documented preprocessing trade-offs and compatibility for calibration, sensitivity, and uncertainty workflows, with expert review still required before decision use.
\end{abstract}
